\chapter{Glossary and where to go next}

\begin{longtable}{>{\RaggedRight\arraybackslash}p{0.32\linewidth}>{\RaggedRight\arraybackslash}p{0.58\linewidth}}
\toprule
\textbf{Term} & \textbf{Read it as} \\
\midrule
\endfirsthead
\toprule
\textbf{Term} & \textbf{Read it as} \\
\midrule
\endhead
Source & The original input database being read. \\[5pt]
Target & The separate database that receives the selected smaller copy. \\[5pt]
Table / row / column & A collection of one kind of record / one record / one kind of value in each record. \\[5pt]
Primary key & The declared identifier column or identifier tuple that uniquely identifies a row. \\[5pt]
Foreign key & A declared reference requiring a child value to match a parent record. \\[5pt]
Parent / child & The referenced record / the record that depends on it. \\[5pt]
Schema & The blueprint of tables, columns, keys, types and rules. \\[5pt]
Constraint / index & A database-enforced rule / a structure that helps database lookups. \\[5pt]
Workload & The chosen list of read-only questions used to assess the copy. \\[5pt]
Closure & A chosen record set plus every parent record it needs recursively. \\[5pt]
Incremental cost & How many additional records must be added after already-selected records are counted. \\[5pt]
Budget & The maximum total rows allowed in the output, including parents. \\[5pt]
Seed & The starting value used to repeat pseudo-random choices. \\[5pt]
Witness & Actual matching record identifiers used to support a query condition/join. \\[5pt]
Heuristic & A practical selection rule that does not prove universal optimality. \\[5pt]
Validation & Checks that structure, identifiers, references, values and row limits meet their rules. \\[5pt]
Acceptance & The explicit decision whether measured query quality is sufficient. \\[5pt]
Transaction / commit / rollback & A group of related pending changes / saving them / cancelling them. \\[5pt]
Survival & An eligible question returns a nonempty result on both original and copy. \\[5pt]
Exact retention & The fraction of original typed answer-row occurrences that still match. \\[5pt]
Group coverage & How many original identifiable groups still appear, irrespective of changed numeric totals. \\[5pt]
Fingerprint & A compact check value derived from observed schema/content, used with a baseline. \\[5pt]
Journal & A durable local record of a job, its progress, outcome and created-schema ownership. \\[5pt]
Benchmark & A planned comparison of repeated experiments with stated input/settings. \\[5pt]
API / function signature & A callable interface / its named inputs, defaults and declared return annotation. \\[5pt]
\bottomrule
\end{longtable}

\section{When you are ready for technical detail}

Use the separate Complete Workflow and Project Review Guide as the technical companion. It contains exact function signatures, class fields, resource-limit keys, source hashes and the full tracked-file manifest. Return to this beginner walkthrough whenever a definition needs a record-level example.

\section{Reproduce the teaching example}

The source bundle includes assets/shop-example.sql and assets/shop-workload.sql, plus worked-example-results.json. The setup file creates four tables and inserts 13 fictional rows; use it only in a fresh empty database. The JSON records the deliberately chosen six-row copy and measured outcomes. Those keys demonstrate one valid selection, rather than claim to benchmark an automatic method.

\section{Compile this document}

Upload the complete beginner source ZIP and select main.tex. XeLaTeX is recommended; the preamble also chooses a compatible Type 1 font setup for pdfLaTeX. Keep the included .tex files and assets folder together. Recompiling changes typesetting, not the date or meaning of saved application measurements.

\section{One final check of your understanding}

You understand the central workflow if you can explain why Line 1001 needs four records, why Line 1002 adds only two, why COUNT survives with a changed answer, and which files choose, copy and check the records.